\documentclass[12pt, a4, german]{article}

\usepackage[T1]{fontenc}
\usepackage[latin1]{inputenc}
\usepackage[ngerman]{babel}	 %use german spellings

\usepackage{../mmm}

\usepackage{color}
\definecolor{grau}{rgb}{0.4,0.4,0.4} 
\definecolor{schwarz}{rgb}{0,0,0} 
\definecolor{BLUE}{rgb}{0.,0.,1.}
\definecolor{RED}{rgb} {1.,0.,0.}
\definecolor{GREEN}{rgb} {0.,.7,0.}

\def\xcolorcmd{\colorlet{black}{gray!150}}
\usepackage{xcolor}

\usepackage{url}
\usepackage{graphicx}
\graphicspath{pics}

\setlength{\parindent}{0em}
\setlength{\textwidth}{17cm}

\newcommand{\red}[1]{\textcolor{RED}{#1}}
\newcommand{\blue}[1]{\textcolor{BLUE}{#1}}
\newcommand{\green}[1]{\textcolor{GREEN}{#1}}

\newcounter{problem}
\newcounter{pointsa}
\newcounter{pointsb}
\newcounter{blatt}

\setcounter{blatt}{1}

\newenvironment{problem}[3]{
  \addtocounter{problem}{1}
  \addtocounter{pointsa}{#1}
  \addtocounter{pointsb}{#2}
  \bigskip 
  {    \textbf{Task \theproblem}: #3 \hrulefill 
%\textsc{Inf}: (\rule{6mm}{0pt} / \textbf{}) - \textsc{Kom}: (\rule{6mm}{0pt} 
 \textbf{}
  }\begin{quote}
  }{  
  \end{quote}}

%http://www.math.tu-dresden.de/~rudl/latex/fonts.pdf

\headheight0.0cm
\topmargin-2.5cm
\oddsidemargin0.0cm
\textwidth15.92cm
\textheight25cm

\usepackage{watermark}
\usepackage{tikz}

\newcommand{\changefont}[3]{
\fontfamily{#1} \fontseries{#2} \fontshape{#3} \selectfont}

\begin{document}

\thiswatermark{
%\centering
%\put(320,-160){\begin{tikzpicture}%
 %           \node (0,0) [opacity=0.05]{%
  %           \includegraphics[width=4.5cm]{../unilogo.png}};%
   %      \end{tikzpicture}}
\put(120,-145){\begin{tikzpicture}%
            \node (0,0) [opacity=0.5]{%
             \includegraphics[width=4.5cm]{../uni-siegel.png}};%
         \end{tikzpicture}}
\put(355,-145){\begin{tikzpicture}%
            \node (0,0) [opacity=0.5]{%
             \includegraphics[width=2.9cm]{../unilogo.pdf}};%
         \end{tikzpicture}}
%\put(115,-145){\includegraphics[width=5cm]{../uni-siegel.png}}
}

\begin{center}
\begin{tabular}[t]{rcl}
\begin{minipage}{3cm}
\vspace*{0.1cm}
\hspace*{-1.8cm}
\includegraphics[width=4.5cm]{../hcm-logo.pdf}
\end{minipage}
&
\begin{minipage}{10cm}
\vspace*{0.5cm}
%\usefont{T1}{phv}{c}{n}\selectfont
\usefont{T1}{phv}{m}{n}\selectfont
{%\normalsize{Universit�t Augsburg}\\
\textcolor{schwarz}{
\normalsize{Institut f�r Informatik\\
Lehrstuhl Human Centered Multimedia\\
\\
\textbf{Prof. Dr. Elisabeth Andr\'e}\\
Simon Flutura
}
}
}
\end{minipage}
\end{tabular}\\
\end{center}
\vspace*{-0.5cm}
\textcolor{schwarz}{\rule{\textwidth}{0.3pt}}
\begin{center}
\usefont{T1}{phv}{b}{n}\selectfont
%\usefont{T1}{uhv}{b}{n}\selectfont
{\Large {\bf{
\usefont{T1}{phv}{b}{n}\selectfont
\vspace*{0.5cm}
\textcolor{schwarz}{SSI on Linux, SS 2015}}}}          
\end{center}
\vspace{0.2cm}
%
%
% =====================================================================
%
\shorthandoff{"} % Use Gaensefuesschen within the text without quoting them
\begin{center}
\usefont{T1}{phv}{b}{n}\selectfont{\large
%\usefont{T1}{uhv}{b}{n}\selectfont{\large
\theblatt. Introduction}
\end{center}
\vspace{0.2cm}
%\changefont{ppl}{m}{n}
%\changefont{ptm}{m}{n}
%\changefont{phv}{m}{n}

%\changefont{cmr}{m}{n}

\changefont{phv}{m}{n}
\normalsize

SSI - Social Signal Iterpretation Framework has been portet to Linux with an Android-port in the works.
Devices like the Raspberry PI, as well as several mobile devices can now be used. This guide shows you how to set up SSI on Linux\\\\
%------------------------------------------------------------------------------------------------------------------------------------------------------------------------------------------------------------
\begin{problem}{0}{0}{Prepare the setup}
Some requirements have to be met to be able to use SSI on Linux.
On any reasonable Distribution you should simply be able to install the dependencies via package manager.


\begin{enumerate}
\item GCC and G++ 4.8 or later have to be installed, because SSI relies on C++11 standard.
\item CMake is needed as build-system.
\item Optionally for graphics output SDL2 development packages are needed as well as cairo development packages.
\item Portaudio development library is needed for audio.
\end{enumerate}
\end{problem}
%------------------------------------------------------------------------------------------------------------------------------------------------------------------------------------------------------------
\begin{problem}{0}{0}{Get SSI}
SSI can be found here:
\begin{enumerate}
	\item from git \verb|https://hcm-lab.de/git/groups/ssi|
\item from subversion \verb|https://hcm-lab.de/svn/Johannes/openssi/trunk/|


\end{enumerate}
\end{problem}

%----------------------------------------------------------------------------------------------------------------------------------------------------------------------------------------
\begin{problem}{0}{0}{Now you are ready to build SSI}


\begin{enumerate}
	\item create a build directory next to the ssi root directory:
	
	\verb|mkdir ./trunk/../build|
	\item run cmake from the build directory to create the build environment:
	\verb|cd build|
	\verb|cmake ../trunk|
	\item run make to build and install ssi(with 6 threads):
	
	\verb|make -j6 install|
\end{enumerate}

to disable GUI use
\verb|SET(SSI_GUI none)|
instead of
\verb|SET(SSI_GUI SDL)|
in \verb|trunk/CMakeLists.txt|

\end{problem}
%----------------------------------------------------------------------------------------------------------------------------------------------------------------------------------------

\begin{problem}{0}{0}{Run tests}
Your freshly build binarys are located in ssi roots \verb|bin_cmake| directory:

\verb|trunk/cmake_bin/Linux|
\begin{enumerate}
	\item To test an audio pipeline prepare and test your microphone.
	then run \verb|./ssiemovoice_test| from \verb|trunk/cmake_bin/Linux|.
	The correct working directory is needed, so the shared plugins can be found.
	\item Apart from c++ code one can as well run xml pipelines via:
	
	\verb|./xmlpipe test.pipeline|


\end{enumerate}
\end{problem}
%----------------------------------------------------------------------------------------------------------------------------------------------------------------------------------------

\begin{problem}{0}{0}{Choose IDE}
To further develop SSI on linux QtCreator is recommended.
Do not forget to install the GNU debugger GDB.

Have fun.

\end{problem}

\end{document}


